\section{Discusión}
\label{sec:discusion}

Los resultados obtenidos permiten analizar el comportamiento del agente desde tres perspectivas: sensibilidad a los hiperparámetros, estabilidad del aprendizaje e interpretación de la política aprendida.

\subsection{Sensibilidad a los hiperparámetros}

El barrido exhaustivo muestra que el rendimiento del agente depende fuertemente de la combinación entre el factor de descuento, la tasa de aprendizaje y el calendario de exploración. La mejor configuración alcanzó un promedio de $431.21$ líneas eliminadas, mientras que otras configuraciones obtuvieron desempeños considerablemente menores. Esta diferencia evidencia que el entrenamiento de agentes de Aprendizaje por Refuerzo no depende únicamente del algoritmo, sino también de una calibración adecuada de sus hiperparámetros.

El factor de descuento fue uno de los componentes más relevantes. Los valores $\gamma=0.990$ y $\gamma=0.999$ presentaron un desempeño promedio muy superior a $\gamma=0.900$. Esto es coherente con el problema de Tetris, donde muchas decisiones correctas no producen una recompensa inmediata, sino que preparan el tablero para obtener mejores recompensas futuras. Por ejemplo, mantener una estructura baja y regular puede no eliminar líneas en el turno actual, pero aumenta la probabilidad de sobrevivir más tiempo y completar más líneas posteriormente.

La tasa de aprendizaje también mostró un efecto importante. Aunque el mejor resultado puntual se obtuvo con $\alpha=0.001$, el promedio agregado de las tasas evaluadas indica que el desempeño no puede explicarse por este parámetro de manera aislada. Una tasa mayor puede acelerar la adaptación de los pesos, pero también aumenta el riesgo de inestabilidad si se combina con configuraciones desfavorables. Por ello, el mejor resultado surge de una interacción favorable entre $\alpha=0.001$, $\gamma=0.999$ y $N_{\varepsilon}=1000$.

\subsection{Exploración y explotación}

El parámetro $N_{\varepsilon}$ controla la rapidez con la que disminuye la exploración del agente. En la mejor configuración, el valor $N_{\varepsilon}=1000$ permitió una transición relativamente temprana hacia la explotación de la política aprendida. Esto sugiere que, para el entorno estudiado, el agente logró adquirir información suficiente durante las primeras etapas del entrenamiento y se benefició de consolidar rápidamente las acciones de mayor valor estimado.

Sin embargo, el análisis agregado muestra que $N_{\varepsilon}=1800$ obtuvo el promedio más alto entre los tres valores evaluados. Esto indica que una exploración más prolongada puede mejorar la robustez general del entrenamiento, aunque no necesariamente produce la mejor configuración individual. Por tanto, existe un compromiso entre explorar durante más tiempo para cubrir mejor el espacio de estados y explotar más temprano para estabilizar la política.

\subsection{Interpretación de la curva de aprendizaje}

La Figura~\ref{fig:curva_aprendizaje} permite observar la evolución del agente durante el entrenamiento. En este tipo de problemas es esperable que la curva no sea estrictamente monótona, debido a la aleatoriedad en la generación de piezas y a la exploración inducida por la política $\varepsilon$-greedy.

Aun con dicha variabilidad, la curva resulta útil para verificar que el agente experimenta una mejora progresiva y que el entrenamiento no corresponde a una ejecución aleatoria sin aprendizaje. La existencia de una tendencia de mejora y posterior estabilización respalda que la función de valor aproxima gradualmente una valoración útil de los estados del tablero.

\subsection{Interpretación de la evolución de pesos}

La Figura~\ref{fig:evolucion_pesos} permite analizar cómo cambian los parámetros de la función de valor durante el entrenamiento. En una aproximación lineal, cada peso determina la influencia de una característica sobre la evaluación del tablero. Por ello, la evolución de los pesos proporciona una forma de interpretación del comportamiento aprendido por el agente.

Una estabilización progresiva de los pesos sugiere que el agente alcanza una valoración más consistente de las características del tablero. En términos prácticos, esto significa que el agente aprende a penalizar configuraciones desfavorables, como tableros altos, huecos internos o superficies irregulares, y a favorecer configuraciones que incrementan la supervivencia y la eliminación de líneas.

\subsection{Alcance del resultado obtenido}

El mejor modelo logró eliminar en promedio $431.21$ líneas antes del estado terminal. Este resultado es relevante porque fue obtenido mediante una aproximación lineal relativamente simple, sin redes neuronales profundas ni búsqueda exhaustiva de largo horizonte. Por tanto, el proyecto demuestra que una representación adecuada del estado mediante características puede ser suficiente para construir un agente efectivo en una versión simplificada del problema de Tetris.

No obstante, el resultado debe interpretarse dentro de las condiciones del entorno implementado. El agente decide por colocaciones completas y no modela maniobras avanzadas dependientes del tiempo real, como desplazamientos durante la caída, rotaciones tardías o técnicas especializadas. En consecuencia, el desempeño obtenido corresponde al modelo de decisión definido en el MDP del proyecto y no necesariamente a todas las variantes posibles del juego Tetris.

\subsection{Comparación con trabajos previos}
El desempeño obtenido, de aproximadamente 431 líneas por partida, se encuentra por debajo de los mejores agentes reportados en la literatura, algunos de los cuales superan decenas de miles de líneas utilizando representaciones más complejas y algoritmos de optimización avanzados. Sin embargo, el resultado obtenido es significativo considerando que el agente utiliza una aproximación lineal y un conjunto reducido de características.

\subsection{Limitaciones}

La principal limitación del enfoque es el uso de una función de valor lineal. Aunque esta elección mejora la interpretabilidad y reduce el costo computacional, también limita la capacidad del agente para capturar relaciones no lineales complejas entre las características del tablero. Además, el entrenamiento sigue siendo sensible a la elección de hiperparámetros, como se evidenció en el barrido experimental.

Otra limitación es que la recompensa se basa principalmente en líneas eliminadas. Aunque esta métrica está alineada con el objetivo del juego, puede no capturar completamente criterios estratégicos intermedios, como preparar el tablero para futuras piezas, reducir riesgos de top-out o mantener opciones abiertas para distintas secuencias de piezas.

\subsection{Implicancias para trabajos futuros}

Los resultados obtenidos abren la posibilidad de extender el sistema hacia enfoques más expresivos, como Deep Q-Learning, redes neuronales convolucionales para representar directamente el tablero o métodos híbridos que combinen heurísticas con Aprendizaje por Refuerzo. También sería posible ampliar el espacio de acciones para incluir maniobras más cercanas al juego real, así como comparar el agente entrenado contra políticas aleatorias, heurísticas clásicas y controladores basados en búsqueda.

En síntesis, el barrido experimental confirma que el agente TD con aproximación lineal es una solución viable y efectiva para el entorno propuesto. La combinación de resultados cuantitativos, curva de aprendizaje y evolución de pesos proporciona evidencia suficiente para sostener que el agente aprende una política útil y no se limita a una estrategia aleatoria o fija.
